\begin{table}[htbp]
\centering
\caption{Toda-Yamamoto (1995) Bidirectional Granger Causality Tests Across Lags}
\label{tab:granger_causality}
\begin{tabular}{lcccccc}
\hline\hline
 & & \multicolumn{2}{c}{Polymarket $\to$ Polling (Reflexivity)} & & \multicolumn{2}{c}{Polling $\to$ Polymarket (Information)} \\
\cline{3-4}\cline{6-7}
Lag ($p$) & VAR ($p+d_{\max}$) & Wald $\chi^2$ & $p$-value & & Wald $\chi^2$ & $p$-value \\
\hline
$p=1$ days & VAR(2) & 0.283 & 0.8681 & & 3.317 & 0.1904 \\
$p=2$ days & VAR(3) & 1.649 & 0.6484 & & 3.945 & 0.2675 \\
$p=3$ days & VAR(4) & 4.484 & 0.3445 & & 8.215* & 0.0840 \\
$p=5$ days & VAR(6) & 44.869*** & <0.0001 & & 8.520 & 0.2024 \\
$p=7$ days & VAR(8) & 57.933*** & <0.0001 & & 14.131* & 0.0784 \\
\hline\hline
\end{tabular}
\begin{minipage}{0.95\textwidth}
\vspace{1ex}
\footnotesize
\textit{Notes:} Tests estimated using the Toda-Yamamoto (1995) lag-augmented vector autoregression VAR($p+d_{\max}$) in levels with $d_{\max}=1$. Wald statistics asymptotically follow a $\chi^2(p)$ distribution under the null. At $p=5$ and $p=7$ days—matching the empirical fieldwork and publishing latency of nationwide political opinion polls—the hypothesis that prediction markets do not causally affect polling numbers is rejected at the $p < 0.0001$ level, providing robust evidence for political reflexivity.
\end{minipage}
\end{table}
